\documentclass[11pt]{article}
\usepackage[margin=1in]{geometry}
\usepackage{amsmath,amssymb,amsthm,mathtools}
\usepackage{hyperref}

\newtheorem{theorem}{Theorem}
\newtheorem{lemma}[theorem]{Lemma}
\newcommand{\N}{\mathbb N}
\newcommand{\Hh}{\mathcal H}

\title{A Quasiperiodic Product of Countably Many\\
Orthogonal Projections Need Not Converge in Norm}
\author{Automatic mathematics run \texttt{fa\_banach\_001}}
\date{11 August 2026}

\begin{document}
\setlength{\emergencystretch}{3em}
\maketitle

\begin{center}
\fbox{\parbox{0.91\textwidth}{
\textbf{Status.} Candidate counterexample, likely valid, for human review.
It gives a negative answer to the countably infinite quasiperiodic-family
question recorded in arXiv:1809.05858.
}}
\end{center}

\section{Source question and answer}

Let $(M_j)$ be closed subspaces of a Hilbert space, let $P_j$ be their
orthogonal projections, and set
\[
 x_n=P_{\sigma(n)}x_{n-1}.
\]
For a finite family, Sakai proved norm convergence when $\sigma$ is
\emph{quasiperiodic}: every index occurs infinitely often and the gaps
between successive occurrences of each fixed index are bounded.

On printed page 16 of Ginat's dissertation \cite{source}, the following
question is attributed to Sakai:
\begin{quote}
Does norm convergence still hold when the family is countably infinite and
the index sequence is quasiperiodic?
\end{quote}

\begin{theorem}\label{thm:main}
There are a countable family $(M_j)_{j\geq1}$ of closed subspaces of
$\ell_2$, a quasiperiodic map $\sigma:\N\to\N$, and $x_0\in\ell_2$ such
that the sequence
\[
 x_n=P_{\sigma(n)}\cdots P_{\sigma(1)}x_0
\]
does not converge in norm.  In fact, $x_n\rightharpoonup0$,
$\lim_n\|x_n\|>0$, and $\bigcap_jM_j=\{0\}$.
\end{theorem}

\section{Idea of the proof}

Move a unit vector from $e_m$ to $e_{m+1}$ in very many tiny planar
rotations.  The sum of the squared angles is finite, so projecting through
the hyperplane implementing each rotation loses only a positive finite
factor.  Nevertheless, the unit directions escape to infinity and converge
weakly to zero.

Each hyperplane is initially used exactly when its rotation is needed.  It is
then withheld until the path has left the two coordinates supporting its
normal vector.  After that time it may be repeated forever without changing
the orbit.  Exact $2$-adic valuation classes provide a disjoint bounded-gap
recurrence schedule for all the old hyperplanes simultaneously.  This makes
the complete index sequence quasiperiodic while preserving the one-pass
nonconvergent orbit exactly.

\section{The slowly rotating path}

Let $(e_m)_{m\geq1}$ be the canonical orthonormal basis of
$\Hh=\ell_2$.  Put
\[
 N_m=2^m,
 \qquad
 \theta_m=\frac{\pi}{2N_m}.
\]
Starting at $v_0=e_1$, join $e_m$ to $e_{m+1}$ by the $N_m$ unit vectors
at successive angles $\theta_m,2\theta_m,\ldots,N_m\theta_m$:
\[
 \cos(r\theta_m)e_m+\sin(r\theta_m)e_{m+1},
 \qquad 1\leq r\leq N_m.
\]
Concatenating these finite strings defines $(v_k)_{k\geq0}$.  Every
$v_k$ is a unit vector, and consecutive vectors lying in stage $m$ satisfy
\[
 \langle v_{k-1},v_k\rangle=\cos\theta_m=:c_k>0.
\]
Moreover $v_k\rightharpoonup0$, since the two coordinates supporting $v_k$
tend to infinity.  At the end of stage $m$ the path equals $e_{m+1}$, so it
does not converge in norm.

For every $k\geq1$, define
\[
 u_k=\frac{v_{k-1}-c_kv_k}{\sqrt{1-c_k^2}},
 \qquad
 Q_k=I-u_k\otimes u_k.
\]
Thus $u_k$ is a unit vector orthogonal to $v_k$, $Q_k$ is the orthogonal
projection onto $u_k^\perp$, and
\begin{equation}\label{eq:rotation}
 Q_kv_{k-1}=c_kv_k.
\end{equation}
If $k$ belongs to stage $m$, then $u_k\in[e_m,e_{m+1}]$.  Consequently
there is an integer $L_k\geq k$ (for example, the end of stage $m+1$) such
that
\begin{equation}\label{eq:safe}
 v_j\perp u_k\quad\text{for every }j\geq L_k.
\end{equation}

Let $a_0=1$ and $a_k=\prod_{i=1}^kc_i$.  The infinite product has a
strictly positive limit.  Indeed,
\[
 \sum_{k\geq1}(1-c_k)
 =\sum_{m\geq1}N_m(1-\cos\theta_m)
 \leq \frac12\sum_{m\geq1}N_m\theta_m^2
 =\frac{\pi^2}{8}\sum_{m\geq1}\frac1{N_m}<\infty.
\]
Since every $c_k>0$, this gives
\begin{equation}\label{eq:positive}
 a_k\longrightarrow a_\infty>0.
\end{equation}

\section{A bounded-gap recurrence schedule}

Write $\nu_2(n)$ for the exact exponent of $2$ dividing the positive
integer $n$.  Let
\[
 P_1=I,\qquad P_{k+1}=Q_k\quad(k\geq1),
\]
and define $\sigma:\N\to\N$ as follows:
\begin{enumerate}
\item $\sigma(2k-1)=k+1$ for $k\geq1$;
\item if $n$ is even, $\nu_2(n)=k+1$ for some $k\geq1$, and
      $n\geq2L_k$, set $\sigma(n)=k+1$;
\item at every remaining even integer, set $\sigma(n)=1$.
\end{enumerate}

This sequence is quasiperiodic.  Index $1$ occurs at least at every integer
$n\equiv2\pmod4$, so its gaps are at most $4$.  For fixed $k\geq1$, index
$k+1$ first occurs at $2k-1$.  After a finite delay it occurs at every
sufficiently large integer satisfying $\nu_2(n)=k+1$.  Those integers form
an arithmetic progression with common difference $2^{k+2}$.  Hence the gap
between successive occurrences of $k+1$ is bounded (the one initial gap is
finite as well).

\section{Exact orbit and failure of norm convergence}

We claim that for every $k\geq1$,
\begin{equation}\label{eq:orbit}
 x_{2k-1}=x_{2k}=a_kv_k.
\end{equation}
This follows by induction.  At odd time $2k-1$, the prescribed projection is
$Q_k$, so \eqref{eq:rotation} sends $a_{k-1}v_{k-1}$ to $a_kv_k$.
At even time $2k$, either the projection is $I$, or it is some $Q_i$ whose
recurrence has been released.  In the latter case $2k\geq2L_i$, hence
$k\geq L_i$, and \eqref{eq:safe} gives $Q_iv_k=v_k$.  Thus the even step
does nothing, proving \eqref{eq:orbit}.

Equations \eqref{eq:positive} and \eqref{eq:orbit} show that
$\|x_n\|\to a_\infty>0$, while $v_k\rightharpoonup0$ gives
$x_n\rightharpoonup0$.  Therefore $(x_n)$ cannot converge in norm.

Finally, the common intersection is zero.  During stage $m$ there are at
least two distinct normal vectors $u_k$ in $[e_m,e_{m+1}]$, and they span
that two-dimensional coordinate plane.  A vector belonging to every
$u_k^\perp$ is therefore orthogonal to $e_m$ and $e_{m+1}$ for every $m$,
so it is zero.  This proves Theorem~\ref{thm:main}.

\section{Bounded literature audit}

The source itself records the countably infinite quasiperiodic case as open.
A later paper of Eskandari and Moslehian \cite{later} proves convergence in a
different \emph{pseudo-periodic} regime: only a finite recurrent core is used,
and its countable tail is subject to additional structure.  That result does
not apply here, because every one of the countably many hyperplane indices is
recurrent with bounded gaps.  Exact-title and exact-phrase searches of arXiv
and the run indexes found no later source containing the counterexample above.
This is a bounded novelty check, not a claim of exhaustive priority.

\begin{thebibliography}{9}
\bibitem{source}
O. Ginat,
\emph{The Method of Alternating Projections},
arXiv:1809.05858 (2018), especially Section 3.4.3, printed p.~16.

\bibitem{later}
R. Eskandari and M. S. Moslehian,
\emph{Convergence of Random Products of Countably Infinitely Many
Projections}, arXiv:2405.04848 (2024).
\end{thebibliography}

\end{document}
